\subsection{Persistance régionale et unicité du sélecteur d’action}
\label{sec:accion-selector-regional}

La microfibre de clôture de la section~\ref{sec:generacion} fournit
l’antécédent régional du lecteur d’action. Avec la réduction ternaire orientée
\(\Psi:\mathcal U_6\to\mathbb F_3^6\) déjà construite, soit
\(\mathcal F_\pi=\Psi^{-1}(010211)\). Ses cinq émissions, séparées
en bloc initial et terminal, sont
\[
\begin{aligned}
 &(501,614,498\mid169,272,272),\\
 &(810,923,870\mid169,272,272),\\
 &(810,983,810\mid169,272,272),\\
 &(870,923,810\mid169,272,272),\\
 &(870,923,810\mid418,674,521).
\end{aligned}
\]
La projection de retour est
\[
 p_{\rm ret}:\mathbb Z^6\longrightarrow\mathbb Z^3,
 \qquad p_{\rm ret}(u_1,\ldots,u_6)=(u_4,u_5,u_6).
\]
Dans cette étape, chaque émission compte une fois : on sélectionne la persistance
entre les cinq régions, en conservant séparément les multiplicités
des germes qui sont à l’origine de ces émissions.

\begin{proposition}[Bloc terminal persistant]
\label{prop:accion-bloque-persistente}
Le multiensemble \(p_{\rm ret}(\mathcal F_\pi)\) a un unique mode,
\[
 b_\pi=(169,272,272),\qquad \operatorname{mult}(b_\pi)=4.
\]
Le bloc restant, \(b_\pi^-=(418,674,521)\), a pour multiplicité un.
\end{proposition}
\begin{proof}
La projection des quatre premières émissions est \(b_\pi\) ; celle de la
cinquième est \(b_\pi^-\). Les deux blocs sont distincts. La multiplicité
maximale est donc quatre et elle est atteinte en un seul bloc.
\end{proof}

Le sélecteur est défini sur les multiensembles de blocs terminaux ayant un
mode unique de type \((r,s,s)\), \(r\ne s\). Le groupe
\(\mathfrak S_3\) agit en renumérotant les positions de ces blocs ;
cette équivariance compare leurs représentations, sans ajouter de symétrie
dynamique à la fibre fixe \(\Psi^{-1}(010211)\).
Le stabilisateur de \(b_\pi\) échange les positions occupées par
\(272\) et fixe la position occupée par \(169\). Pour un bloc de type
\((r,s,s)\), avec \(r\ne s\), la valeur de cette position fixe est
\[
 \operatorname{sing}(b)=\sum_{j=1}^3 b_j-2\operatorname{moda}(b).
\]
La sélection du mode régional suivie de cette opération définit le
sélecteur résiduel \(\rho_\pi\).

\begin{theorem}[Unicité équivariante du retour]
\label{thm:accion-selector-169}
Parmi les sélecteurs invariants par permutation des cinq régions,
équivariants par permutation des positions terminales, qui
sélectionnent le bloc de multiplicité maximale et retiennent une position
fixée par son stabilisateur, la valeur de retour est unique :
\[
 \rho_\pi=\operatorname{sing}(b_\pi)=169.
\]
\end{theorem}
\begin{proof}
L’invariance par rapport aux régions fait dépendre la sélection du
multiensemble, non de l’ordre d’énumération. La persistance modale
sélectionne \(b_\pi\) par la proposition~\ref{prop:accion-bloque-persistente}.
Son stabilisateur a deux orbites de positions : une paire et un point fixe.
La condition de position fixe sélectionne ce dernier, dont la valeur est
\(169\). Lorsqu’on permute les positions, l’équivariance transporte cette
position et conserve sa valeur. La construction décrite satisfait les
quatre conditions, ce qui prouve également l’existence.
\end{proof}

La différence orientée entre les deux retours conserve une donnée supplémentaire :
\[
 \omega_\pi=b_\pi^- - b_\pi=(249,402,249),\qquad
 \langle(1,1,1),\omega_\pi\rangle=900.
\]
Le sélecteur \(169\) est une donnée terminale de degré zéro. La différence
\(\omega_\pi\) assigne une donnée de degré un à l’arête orientée entre
les deux blocs ; son traitement comme cocycle exige de spécifier le complexe
et de vérifier le cobord correspondant. Cette distinction conserve tant
le retour que sa variation, sans les utiliser indifféremment.

La sélection est régionale : dans le catalogue complet de \(468\) émissions,
\(169\) apparaît \(84\) fois, \(14\) dans chaque position. Sa fonction dans
le lecteur provient conjointement de la microfibre, de la persistance modale et du
stabilisateur ; la fréquence globale est conservée comme un autre
dénombrement du même catalogue.

\subsection{Longueur de mot et caractère de retour défini par un déterminant}
\label{sec:accion-grado-determinante}

Soit \(\mathscr P_{\rm pal}=\bigoplus_{n\ge0}\mathscr P_{{\rm pal},n}\) le module libre
sur les mots finis de l’alphabet d’émission, gradué par la longueur,
avec la filtration induite. Ce domaine contient les émissions régionales
et leurs concaténations ; la concaténation additionne les degrés. Le caractère
formel de longueur
\[
 \chi_z([w])=z^{|w|},\qquad
 \chi_z([uv])=\chi_z([u])\chi_z([v])
\]
conserve cette opération. La compatibilité cylindrique du lecteur envoie
\(\mathscr P_{{\rm pal},n}\) sur la composante homogène \(z^n\mathbb R\) de
\(\mathbb R[z]\) ; l’évaluation \(z=\alpha\) s’effectue ensuite.
Comme chaque \(w\in\mathcal F_\pi\) a une longueur de six, l’impulsion
résiduelle est indépendante de la région choisie et de degré six :
\[
 \rho_\pi\chi_z([w])=169z^6,
 \qquad\left.\rho_\pi\chi_z([w])\right|_{z=\alpha}=169\alpha^6.
\]
La sélection enregistre une impulsion commune ; elle ne somme pas de nouveau les cinq
émissions après en avoir extrait le retour persistant.
La longueur de l’émission et celle d’une marche dans le graphe
d’incidence sont des degrés d’objets différents. En particulier, les marches
fermées qui parcourent les quatre ancrages régionaux ont une longueur
minimale de huit ; le degré utilisé ici compte les six positions d’une
émission, non les arêtes de cette marche.

La clôture dodécaphasique et la complétion des trois paires nonadiques
fournissent la coordonnée commune \(A=10^3\alpha\). En radians,
\[
 x=\frac{\pi A}{180}=\frac{50\pi}{9}\alpha.
\]
Soit \(R_{\rm h}\) une involution réelle bidimensionnelle de trace nulle,
\(R_{\rm h}^2=I_2\), \(\operatorname{tr}R_{\rm h}=0\). Le transport bicouche est
\[
 \mathcal T(x,y)=\exp(-xI_2-yR_{\rm h}).
\]
La coordonnée orientée \(y\) reste non évaluée. Comme le polynôme
\(X^2-1\) a des racines distinctes et la trace de \(R_{\rm h}\) est nulle,
ses valeurs propres sont \(1,-1\). Dans une base de vecteurs propres, les deux
facteurs du transport sont \(e^{-x-y}\) et \(e^{-x+y}\). Par conséquent,
\[
 \det\mathcal T=e^{-x-y}e^{-x+y}=e^{-2x}
 =\exp\!\left(-\frac{100\pi\alpha}{9}\right)=D_A.
\]
Sur la droite déterminante, \(\bigwedge^2\mathcal T\) agit comme
\(D_A\) fois l’identité. Le produit élimine la coordonnée orientée et conserve la contraction
commune. Le déterminant est invariant par conjugaison et par
l’échange de feuillets \(R_{\rm h}\mapsto-R_{\rm h}\) ; on construit ainsi \(D_A\)
avant d’évaluer le contre-angle.

\begin{proposition}[Caractère multiplicatif du retour]
\label{prop:accion-caracter-determinantal}
Le caractère multiplicatif \(\delta_A:(\mathbb N_0,+)\to(\mathbb R_{>0},\cdot)\)
dont la valeur sur un prolongement élémentaire est \(D_A\) est déterminé
par \(\delta_A(k)=D_A^k\).
\end{proposition}
\begin{proof}
Un caractère préserve l’unité, donc \(\delta_A(0)=1\).
La concaténation de \(k\) prolongements donne
\(\delta_A(k)=\delta_A(1)^k=D_A^k\). Ces puissances satisfont
la multiplicativité et la valeur élémentaire prescrite, prouvant l’existence
et l’unicité.
\end{proof}

\Needspace{17\baselineskip}
\begin{definition}[Règles de la continuation régionale]
\label{def:accion-lector-regional}
Le lecteur régional défini par le déterminant du transport conserve cinq opérations :
\begin{enumerate}
 \item la longueur cylindrique détermine le degré analytique ;
 \item l’impulsion de retour et la continuation stricte sont des termes
 distincts d’un renouvellement ;
 \item après avoir enregistré \(\rho_\pi\) une seule fois, la continuation
 commence à l’unité de la droite déterminante ;
 \item chaque prolongement agit par \(\bigwedge^2\mathcal T\), dont
 le facteur scalaire est \(D_A\) ;
 \item dans l’orientation cardinale d’APP, le retour régional occupe le
 secteur compensateur et se soustrait à la phase d’action de cinquième degré.
\end{enumerate}
\end{definition}

La normalisation unitaire détermine l’amplitude de la continuation, outre
sa raison. En effet, toutes les séries
\[
 F_\lambda(z)=169z^6+\lambda\frac{D_Az^7}{1-D_Az}
 \qquad(\lambda\in\mathbb R)
\]
conservent l’impulsion finie et la même récurrence multiplicative des
coefficients de la queue. La troisième règle fixe \(\lambda=1\) avant d’évaluer
\(z=\alpha\). Le résidu s’enregistre une fois ; le prolongement compte
la droite qui se poursuit, avec son unité. L’équation de renouvellement
\eqref{eq:revision-renovacion} et sa preuve d’unicité, développées
ci-après, complètent cette chaîne entre région, degré,
transport et section d’action.
